\begin{thebibliography}{99}
\bibitem[1]{ref01} I. Belegradek \& D. V. Osin - “Rips construction and Kazhdan property (T)” , Groups Geom. Dyn. 2 (2008) no. 1, p. 1-12
\bibitem[2]{ref02} B. H. Bowditch - “Tight geodesics in the curve complex” , Invent. Math. 171 (2008) no. 2, p. 281-300
\bibitem[3]{ref03} M. R. Bridson \& A. Haefliger - Metric spaces of non-positive curvature , Grundlehren Math. Wiss. , vol. 319 , Springer-Verlag, Berlin, 1999
\bibitem[4]{ref04} M. Coornaert, T. Delzant \& A. Papadopoulos - Géométrie et théorie des groupes. Les groupes hyperboliques de Gromov , Lect. Notes in Math. , vol. 1441 , Springer-Verlag, Berlin, 1990
\bibitem[5]{ref05} R. Coulon - “Asphericity and small cancellation theory for rotation families of groups” , Groups Geom. Dyn. 5 (2011) no. 4, p. 729-765
\bibitem[6]{ref06} R. Coulon - “On the geometry of Burnside quotients of torsion free hyperbolic groups” , Internat. J. Algebra Comput. 24 (2014) no. 3, p. 251-345
\bibitem[7]{ref07} R. Coulon - “Partial periodic quotients of groups acting on a hyperbolic space” , Ann. Inst. Fourier (Grenoble) 66 (2016) no. 5, p. 1773-1857
\bibitem[8]{ref08} F. Dahmani, V. Guirardel \& D. V. Osin - Hyperbolically embedded subgroups and rotating families in groups acting on hyperbolic spaces , Mem. Amer. Math. Soc. , vol. 245, no. 1156 , American Mathematical Society, Providence, RI, 2017
\bibitem[9]{ref09} T. Delzant \& M. Gromov - “Courbure mésoscopique et théorie de la toute petite simplification” , J. Topology 1 (2008) no. 4, p. 804-836
\bibitem[10]{ref10} É. Ghys \& P. de la Harpe (eds.) -- Sur les groupes hyperboliques d’après Mikhael Gromov, Progress in Math., vol. 83, Birkhäuser Boston, Inc., Boston, MA, 1990.
\bibitem[11]{ref11} M. Gromov - “Hyperbolic groups” , in Essays in group theory , Springer, New York, 1987, p. 75-263
\bibitem[12]{ref12} M. Gromov - Mesoscopic curvature and hyperbolicity , Contemporary Mathematics , vol. 288 , American Mathematical Society, Providence, RI, 2001
\bibitem[13]{ref13} D. Groves \& J. F. Manning - “Dehn filling in relatively hyperbolic groups” , Israel J. Math. 168 (2008) no. 1, p. 317-429
\bibitem[14]{ref14} A. Martino \& A. Minasyan - “Conjugacy in normal subgroups of hyperbolic groups” , Forum Math. 24 (2012) no. 5, p. 889-910
\bibitem[15]{ref15} A. Y. Ol’shanskii - “On residualing homomorphisms and G-subgroups of hyperbolic groups” , Internat. J. Algebra Comput. 3 (1993) no. 4, p. 365-409
\bibitem[16]{ref16} D. V. Osin - “Acylindrically hyperbolic groups” , Trans. Amer. Math. Soc. 368 (2016) no. 2, p. 851-888
\bibitem[17]{ref17} J.-P. Serre - Arbres, amalgames, $\mathrm{SL}_2$ , Astérisque , vol. 46 , Société Mathématique de France, Paris, 1977
\end{thebibliography}
